\section{Outer estimates used by the anchored construction}
\label{sec:anchor-toolkit-outer-estimates}

This section collects the outer facts used in the anchored connection
argument.  All constants below are independent of the number of vertices and
of the admitted Markov and delay layers.  We use fold time \(s=\delta t\), so
the delays are the fixed translations \(s\mapsto s-\theta_k\).  No backward
RFDE semiflow is used: backward estimates occur only on a transported
one-dimensional unstable line.

Put
\begin{equation}
 \mathcal A_N=D(P_N-\Id),\qquad
 \mathcal N_N(x)=c_N\circ x^{\circ2}+\frac\beta3x^{\circ3}.
 \label{eq:anchor-toolkit-fast-objects}
\end{equation}

\begin{proposition}[Critical curve and frozen history splitting]
\label[proposition]{prop:anchor-toolkit-critical-splitting}
Assume \cref{eq:shared-markov,eq:shared-curvature,eq:shared-layer-bounds}.
There are \(r_0,C>0\), independent of \(N\), such that for \(|r|\le r_0\)
the equation
\begin{equation}
 A_Nz_{0,N}(r)=P_{\perp,N}\mathcal N_N
       \bigl(r\mathbf1+z_{0,N}(r)\bigr),
 \qquad \|z_{0,N}(r)\|_N\le Cr^2,
 \label{eq:anchor-toolkit-critical-transverse}
\end{equation}
has a unique solution in the displayed ball.  With
\begin{equation}
 V_{0,N}(r)=\mathbf1+r\mathbf1+z_{0,N}(r),\qquad
 w_{0,N}(r)=\frac23-\pi_N^T\mathcal N_N
       \bigl(r\mathbf1+z_{0,N}(r)\bigr),
 \label{eq:anchor-toolkit-critical-curve}
\end{equation}
the constant voltage history is an exact zero of the frozen voltage field,
for every positive \(\delta\) and every admitted redistribution.  Moreover,
uniformly through three (r)-derivatives,
\begin{align}
 z_{0,N}(r)&=g_Nr^2+g_{3,N}r^3+O_{\rm unif}(r^4),\notag\\
 w_{0,N}(r)&=\frac23-\alpha r^2-\kappa_Nr^3
                 +O_{\rm unif}(r^4),
 \qquad
 g_{3,N}=2A_N^{-1}P_{\perp,N}(c_N\circ g_N).
 \label{eq:anchor-toolkit-critical-expansion}
\end{align}

The matrix
\begin{equation}
 \mathcal J_N^{\rm fv}(r)=\mathcal A_N-
 \diag\!\left(2c_N\circ\{V_{0,N}(r)-\mathbf1\}
       +\beta\{V_{0,N}(r)-\mathbf1\}^{\circ2}\right)
 \label{eq:anchor-toolkit-frozen-matrix}
\end{equation}
has one simple real eigenvalue \(\lambda_{c,N}(r)\), with right and left
objects \(e_{c,N}(r)\), \(\ell_{c,N}(r)\) normalized by
\(\ell_{c,N}\mathbf1=\ell_{c,N}e_{c,N}=1\).  If
\(P_N^s=\Id-e_{c,N}\ell_{c,N}\), then
\begin{align}
 \lambda_{c,N}(r)&=-2\alpha r-3\kappa_Nr^2+O_{\rm unif}(r^3),\notag\\
 \|e^{\mathcal J_N^{\rm fv}(r)t}P_N^s(r)\|
 &\le Ce^{-D\gamma t/2},\qquad t\ge0,\notag\\
 \lambda_{c,N}(r)&=\frac{w_{0,N}'(r)}{\Theta_N(r)},
 \qquad
 \Theta_N(r):=\ell_{c,N}(r)V_{0,N}'(r),
 \qquad \frac12\le\Theta_N(r)\le\frac32.
 \label{eq:anchor-toolkit-frozen-identities}
\end{align}
Thus the voltage layer is attracting for \(r>0\) and has precisely one
unstable direction for \(r<0\).

Let
\begin{equation}
 \mathscr L_\delta=\log(1/\delta),\qquad
 S_\delta=\sqrt{2p\mathscr L_\delta},\qquad
 \rho_\delta=\frac{\delta S_\delta}{2\alpha},
 \label{eq:anchor-toolkit-fold-scales}
\end{equation}
where \(p>4\) is fixed.  For the shared-resource equation the reduced speed
on \(\rho_\delta\le |r|\le r_0\) is
\begin{equation}
 q_{N,\delta,\nu}^{0}(r)
 =\frac{r-\delta^2\nu}{w_{0,N}'(r)}
 =-\frac1{2\alpha}
   +O_{\rm unif}\!\left(|r|+\frac{\delta^2}{|r|}\right),
 \qquad -q_+\le q^0\le-q_-<0.
 \label{eq:anchor-toolkit-reduced-speed}
\end{equation}
In the fold coordinate \(X=r/\delta\), the frozen critical curve at
\(r=\pm\rho_\delta\) has
\begin{equation}
 \frac{2/3-w_{0,N}(r)}{\delta^2}-\alpha X^2
       +\frac1{2\alpha}
 =\frac1{2\alpha}+O_{\rm unif}(\delta S_\delta^3).
 \label{eq:anchor-toolkit-critical-fold-offset}
\end{equation}

The matrix splitting persists on the frozen-resource voltage history space
\(\mathscr X_N=C([-\theta_m,0];\mathbb R^N)\).  Define
\begin{align}
 \mathfrak a_N(r)&=
 \left\{|\lambda_{c,N}(r)|^{-1}+2(D\gamma)^{-1}\right\}^{-1},\notag\\
 \kappa_{N,\delta}(r)&=\vartheta
 \left\{\frac\delta{\mathfrak a_N(r)}
          +\frac{\theta_m}{\mathscr L_\delta}\right\}^{-1},
 \qquad 0<\vartheta<\frac14,\notag\\
 \|\phi\|_{N,r,\delta}^{\sharp}
 &=\sup_{-\theta_m\le\theta\le0}
       e^{2\kappa_{N,\delta}(r)\theta}\|\phi(\theta)\|_N.
 \label{eq:anchor-toolkit-fading-norm}
\end{align}
For \(\rho_\delta\le|r|\le r_0\) and for the structural layers
\(L_{k,N}(\eta)=B_{k,N}+E_k\), the frozen equation
\begin{equation}
 u'=\delta^{-1}\mathcal J_N^{\rm fv}(r)u
 +\delta K\sum_{k=0}^mL_{k,N}(\eta)
       \{u(s)-u(s-\theta_k)\}
 \label{eq:anchor-toolkit-frozen-rfde}
\end{equation}
is exponentially stable when \(r>0\).  When \(r<0\), it has an invariant
splitting \(\mathscr X_N=\mathscr W^s\oplus\mathscr W^u\), with
\(\operatorname{codim}\mathscr W^s=\dim\mathscr W^u=1\), uniformly bounded
projections, and
\begin{align}
 \|\mathcal T(s)\phi^s\|_{N,r,\delta}^{\sharp}
 &\le Ce^{-\kappa_{N,\delta}(r)s}
       \|\phi^s\|_{N,r,\delta}^{\sharp},\qquad s\ge0,\notag\\
 \|\mathcal T(-s)|_{\mathscr W^u}\|
 &\le Ce^{-\lambda^u_{N,\delta,r,\eta}s},\qquad s\ge0,\notag\\
 \left|\lambda^u_{N,\delta,r,\eta}
       -\delta^{-1}\lambda_{c,N}(r)\right|&\le C\delta.
 \label{eq:anchor-toolkit-frozen-history-bounds}
\end{align}
The splitting and its normalized unstable history have two uniformly
bounded \(\eta\)-derivatives in fixed graph charts.  The second line is only
the inverse on \(\mathscr W^u\), not a backward RFDE solution operator.
\end{proposition}

\begin{proof}
The Dobrushin estimate gives
\(\|A_N^{-1}\|_{E_N\to E_N}\le(D\gamma)^{-1}\).  The right-hand side of
\cref{eq:anchor-toolkit-critical-transverse} is therefore a contraction on
\(\|z\|_N\le Cr^2\), uniformly in \(N\).  Differentiating its fixed-point
equation gives uniform finite (r)-jets.  Comparing powers of (r) and
then projecting with \(\pi_N^T\) proves
\cref{eq:anchor-toolkit-critical-expansion}.  Constant histories annihilate
every delayed difference, which proves the exact frozen-voltage assertion.

In the splitting
\(\mathbb R^N=\operatorname{span}\{\mathbf1\}\oplus E_N\), the matrix
\(\mathcal J_N^{\rm fv}(r)\) is an \(O(|r|)\) perturbation of
\(0\oplus A_N\).  The two invariant graph equations are contractions with
linear inverses \(A_N^{-1}\) and its dual.  They give the simple eigenpair,
the stable semigroup bound, and the expansions above.  Differentiating the
critical equation yields
\[
 \mathcal J_N^{\rm fv}(r)V_{0,N}'(r)=w_{0,N}'(r)\mathbf1;
\]
left multiplication by \(\ell_{c,N}(r)\) proves the exact last identity in
\cref{eq:anchor-toolkit-frozen-identities}.  The speed and fold-offset
formulas follow by substitution in the critical expansion.

For the history equation, use the current stable Green kernel forward and,
on the repelling side, the current unstable kernel backward on its
one-dimensional range.  Prescribed old histories are translated only
forward.  In the norm \cref{eq:anchor-toolkit-fading-norm}, the delayed
perturbation of the current Green map has norm at most
\[
 C\frac{\delta^2}{\mathfrak a_N(r)}
       \{1+e^{2\kappa_{N,\delta}(r)\theta_m}\}
 \le C\frac{\delta^{1-2\vartheta}}{S_\delta}=o(1).
\]
The Neumann series gives the stable and mixed Green operators and
\cref{eq:anchor-toolkit-frozen-history-bounds}.  Differentiating the same
fixed graph equations twice in \(\eta\) proves the asserted uniform jets.
\end{proof}

\begin{proposition}[Exact phase and resource quotient]
\label[proposition]{prop:anchor-toolkit-resource-phase-quotient}
Let \(Z_*(s)=\mathcal H_{V,w,q}(r_s)\) be an exact generated history of the
shared-resource RFDE on a delay-enlarged interval, with
\(r_s'=\delta q(r_s)\), \(q\ne0\), and \(w'(r_s)\ne0\).  If
\(\mathcal U_*(s,t)\) is the full variational process, then
\begin{equation}
 E_s^c:=\partial_r\mathcal H_{V,w,q}(r_s),\qquad
 \mathcal U_*(s,t)E_t^c=\frac{q(r_s)}{q(r_t)}E_s^c.
 \label{eq:anchor-toolkit-center-cocycle}
\end{equation}
Writing \(\varrho_N(\phi,\omega)=\pi_N^T\phi(0)\), the normalized time
tangent and exact phase quotient are
\begin{equation}
 \tau_s=\frac{(\partial_sV_s,\partial_sw_s)}{r_s'},\qquad
 Q_s=\Id-\tau_s\varrho_N,qquad
 \widehat{\mathcal U}_*(s,t)
 =Q_s\mathcal U_*(s,t)|_{\ker\varrho_N}.
 \label{eq:anchor-toolkit-phase-quotient}
\end{equation}
They satisfy \(\varrho_N\tau_s=1\) and the cocycle identity on
\(\ker\varrho_N\).

For a full variation \(\bigl(\xi_s,\omega(s)\bigr)\), put
\begin{align}
 a(s)&=\frac{\omega(s)}{w'(r_s)},\notag\\
 e_s(\theta)&=\frac{q(r_{s+\theta})}{q(r_s)}V'(r_{s+\theta}),\notag\\
 \psi_s(\theta)&=\xi(s+\theta)-e_s(\theta)a(s),
 \qquad u(s)=\psi_s(0).
 \label{eq:anchor-toolkit-resource-coordinates}
\end{align}
Then \(\psi_s\) is unchanged on adding a multiple of \(E_s^c\), and the
quotient equations are exactly
\begin{align}
 a'&=\delta q'(r_s)a+\frac\delta{w'(r_s)}\pi_N^Tu,\notag\\
 (\partial_s-\partial_\theta)\psi_s(\theta)
 &=-\delta\frac{e_s(\theta)}{w'(r_s)}\pi_N^Tu,\notag\\
 u'&=\delta^{-1}\mathcal J_*(s)u
 +\delta K\sum_{k=0}^mL_{k,N}(\eta)
       \{u(s)-\psi_s(-\theta_k)\}
 -\delta\frac{V'(r_s)}{w'(r_s)}\pi_N^Tu.
 \label{eq:anchor-toolkit-resource-system}
\end{align}
On the voltage history space define
\begin{equation}
 (\mathcal K_s\phi)(\theta)=\delta\int_\theta^0
 \frac{e_{s+\sigma}(\theta-\sigma)}{w'(r_{s+\sigma})}
       \pi_N^T\phi(\sigma)\,\dd\sigma.
 \label{eq:anchor-toolkit-volterra-operator}
\end{equation}
Then
\begin{equation}
 \psi_s=(\Id-\mathcal K_s)u_s,
 \label{eq:anchor-toolkit-history-conjugacy}
\end{equation}
and (u) satisfies the ordinary nonautonomous RFDE obtained from the last
line of \cref{eq:anchor-toolkit-resource-system} by replacing
\(\psi_s(-\theta_k)\) with
\(u(s-\theta_k)-(\mathcal K_su_s)(-\theta_k)\).

If \(q,q^{-1},V'\) are uniformly bounded and
\(|w'(r)|\ge c_w|r|\ge c_w c_\rho\rho_\delta\) on the collar, then in every
fading norm of \cref{eq:anchor-toolkit-fading-norm},
\begin{equation}
 \|\mathcal K_s\|\le C\frac\delta{|r_s|}
       \le\frac C{S_\delta},\qquad
 \|(\Id-\mathcal K_s)^{-1}\|\le2.
 \label{eq:anchor-toolkit-volterra-bounds}
\end{equation}
Thus the full anisotropic norm
\(\|\xi_s\|^\sharp+|\omega|/|w'(r_s)|\) is uniformly equivalent to
\(\|\psi_s\|^\sharp+|a|\).
\end{proposition}

\begin{proof}
The derivative of an exact generated history is \(E_s^c\); invariance of
the autonomous time tangent gives \cref{eq:anchor-toolkit-center-cocycle}.
The cocycle identity in \cref{eq:anchor-toolkit-phase-quotient} follows
because the omitted term is transported into the time-tangent line and is
annihilated by \(Q_s\).

The full variational resource row is \(\omega'=\delta\pi_N^T\xi(0)\).
Differentiate \(a=\omega/w'(r_s)\) and use the exact resource identity
\[
 q(r)w'(r)=r-\delta^2\nu,
 \qquad q'w'+qw''=1=\pi_N^TV'.
\]
This gives the first line of \cref{eq:anchor-toolkit-resource-system}.
Subtracting \(aE_s^c\) from the full variational equation gives the other
two lines; all tangent terms cancel.  Integration of the transport equation
along the characteristic
\(\sigma\mapsto(s+\sigma,\theta-\sigma)\) proves
\cref{eq:anchor-toolkit-history-conjugacy} with the sign displayed in
\cref{eq:anchor-toolkit-volterra-operator}.  Finally,
\(\|e_s/w'(r_s)\|\le C/|r_s|\); the weighted integral over one fixed delay
gives the first bound in \cref{eq:anchor-toolkit-volterra-bounds}, and the
second follows from its Neumann series.
\end{proof}

\begin{proposition}[Canonical exact outer reference and coherent hits]
\label[proposition]{prop:anchor-toolkit-reference-hits}
Fix a compact structural parameter box
\(p=(\nu,\eta)\in\mathcal P_N\), a branch
\(\sigma\in\{a,r\}\), and a compatible raw-collar chart of order at least
eight at the near-fold section.  Put
\begin{equation}
 \mathfrak b_{\delta,\varkappa}
 =\varkappa^2+S_\delta^{-2}+\sqrt\delta\,S_\delta^{7/2},
 \qquad
 \frac1{\varkappa^2S_\delta^3}=o(1).
 \label{eq:anchor-toolkit-diagonal-wedge}
\end{equation}
After decreasing the parameter box and then \(\delta_0\), the exact
fixed-clock raw residual has a canonical zero \(X_{\sigma,p}\) in a ball of
radius \(C\mathfrak b_{\delta,\varkappa}\) about its affine predictor.
The zero reconstructs a classical forward segment of the unmodified RFDE;
its old history satisfies the compatibility recursion and all delayed values
after one delay are generated by that same segment.  The result remains true
for compatible raw-collar perturbations of bounded normalized size, but only
the canonical zero is used below.

The nonlinear correction \(h_{\sigma,p}\) satisfies, in the component norm
used by the central--outer overlap,
\begin{equation}
 \|h_{\sigma,p}\|_*
 \le C\mathfrak b_{\delta,\varkappa},
 \qquad
 \|D_p^{\rm sel}h_{\sigma,p}\|_*
 \le C\delta^{-M}\mathfrak b_{\delta,\varkappa}.
 \label{eq:anchor-toolkit-reference-correction}
\end{equation}
On its generated graph the resource component can be written
\begin{equation}
 w_{\sigma,p}(r)=w_{0,N}(r)
   +\delta^2\mathcal D_{\sigma,p}(r)-\delta^2q^0(r)-p_{\sigma,p}(r),
 \label{eq:anchor-toolkit-resource-defect}
\end{equation}
where, on the reserved near-fold collar,
\begin{equation}
 \|\mathcal D_{\sigma,p}\|_{C^1}\le C\delta S_\delta^2,
 \qquad
 \|p_{\sigma,p}\|_{C^1}\le
 C\delta^2\mathfrak b_{\delta,\varkappa}.
 \label{eq:anchor-toolkit-backtrack-budget}
\end{equation}
In particular, for \(r=-c_b\rho_\delta\), with fixed \(c_b>1\),
\begin{equation}
 w_{r,p}(r)=w_{0,N}(r)+\frac{\delta^2}{2\alpha}
       +O(\delta^3S_\delta^2
             +\delta^2\mathfrak b_{\delta,\varkappa}).
 \label{eq:anchor-toolkit-dynamic-overlap}
\end{equation}

Fix \(c_c>0\), and let
\begin{equation}
 \rho_{b,\delta}=c_b\rho_\delta,
 \qquad \rho_{c,\delta}=c_c\delta\mathscr L_\delta.
 \label{eq:anchor-toolkit-core-sections}
\end{equation}
On the repelling reference there are unique hits
\begin{equation}
 r_p(s_b)=-\rho_{b,\delta},\qquad
 r_p(s_c)=-\rho_{c,\delta},\qquad
 (M+2)\theta_m<s_b<s_c<S_p-\theta_m.
 \label{eq:anchor-toolkit-generated-hits}
\end{equation}
Every delay window based between \(s_b\) and \(S_p\) is therefore a literal
generated history.

If \(e_h\) is any fixed regular section crossed by the reference and
\(S_h(p)\) is its hit time, then the event-relative history
\begin{equation}
 \mathcal E_h(p)=
 \left(\theta\mapsto v_p(S_h(p)+\theta),w_p(S_h(p))\right)
 \label{eq:anchor-toolkit-event-hit}
\end{equation}
has the selected jets \(1,\partial_\nu,D_\eta\) in one fixed strong-to-weak
chart.  More generally, for the rectangular parameter jet used in the
anchored proof and \(b=i+j\le3\),
\begin{equation}
 \|\partial_\nu^iD_\eta^j\mathcal E_h(p)\|_{\mathfrak S^3}
 \le C_b\delta^{-2b}.
 \label{eq:anchor-toolkit-event-jets}
\end{equation}
The event derivative is the exact phase projection
\begin{equation}
 D\mathcal E_h(p)\dot p
 =\bigl(\Id-\tau_h\varrho_N\bigr)
       D_pZ_p(S_h(p))\dot p.
 \label{eq:anchor-toolkit-event-derivative}
\end{equation}
\end{proposition}

\begin{proof}
Work on the fixed fold-time interval and its fixed-delay enlargement.  Use
the transverse current variable, the resource defect, and a scalar phase
path as unknowns.  The linearized distributed equation is solved by the
stable current Green operator and, on the repelling branch, by the inverse
on the single current-unstable line.  The resource row is then integrated
forward and the nonzero endpoint speed closes the phase row.  In the
anisotropic norm in which a distributed bulk source, a compatible collar,
and a phase displacement have their natural scales, this triangular inverse
has norm bounded by one constant.

The voltage nonlinearity is cubic, the delay supports are fixed, and the
finite compatibility recursion is polynomial.  Consequently the nonlinear
remainder \(\mathcal Q_{\sigma,p}\), at the affine predictor
\(X^{\rm aff}_{\sigma,p}\), obeys
\begin{align}
 \|\mathcal Q_{\sigma,p}(X^{\rm aff}_{\sigma,p})\|_*^{\rm ran}
 &\le C\mathfrak b_{\delta,\varkappa},\notag\\
 \|\mathcal Q_{\sigma,p}(X^{\rm aff}_{\sigma,p}+h)
       -\mathcal Q_{\sigma,p}(X^{\rm aff}_{\sigma,p}+k)\|_*^{\rm ran}
 &\le C\mathfrak b_{\delta,\varkappa}\|h-k\|_*
 \label{eq:anchor-toolkit-reference-two-point}
\end{align}
on a fixed star ball.  The contraction theorem gives the canonical zero,
its exact reconstruction, and the first estimate in
\cref{eq:anchor-toolkit-reference-correction}.  Differentiating the
selected zero equation in the fixed chart gives the second estimate; it
does not differentiate an RFDE process in operator topology.

For the resource formula, insert the generated backtracks in the exact
resource identity.  Their displacement from the reduced backtracks is
\(O(\delta S_\delta^2)\); row neutrality removes the common component.
Together with \cref{eq:anchor-toolkit-reduced-speed} this proves
\cref{eq:anchor-toolkit-resource-defect,eq:anchor-toolkit-dynamic-overlap}.
The uniform negative speed gives the ordered hits and their delay buffers.
Finally the ordinary implicit-function theorem applied to
\(r_p(S_h)-e_h=0\) gives
\(DS_h=-D_pr_p/r_p'\); substitution into the differentiated history is
exactly \cref{eq:anchor-toolkit-event-derivative}.  The fixed interval has
the spare time derivatives needed for the strong event window; the chain
rule gives the safe loss in \cref{eq:anchor-toolkit-event-jets}.
\end{proof}

\begin{proposition}[Generated-core inverse and action estimates]
\label[proposition]{prop:anchor-toolkit-generated-core}
Retain the repelling reference of
\cref{prop:anchor-toolkit-reference-hits} and the quotient of
\cref{prop:anchor-toolkit-resource-phase-quotient}.  On
\(I_p=[s_b,S_p]\), put
\begin{align}
 a_p(s)&=\mathfrak a_N(r_p(s)),\notag\\
 \kappa_p(s)&=\vartheta_*
 \left\{\frac\delta{a_p(s)}
          +\frac{\theta_m}{\mathscr L_\delta}\right\}^{-1},
 \qquad 0<\vartheta_*<\vartheta,\notag\\
 \mathscr A_p(s,t)&=\int_t^s\kappa_p(u)\,\dd u,\notag\\
 \|\phi\|_{p,s,\sharp}
 &=\sup_{-\theta_m\le\theta\le0}
   e^{-2\mathscr A_p(s,s+\theta)}\|\phi(\theta)\|_N.
 \label{eq:anchor-toolkit-core-action-norm}
\end{align}
Then \(a_p(s)\asymp|r_p(s)|\), \(q_p\) is bounded away from zero, and
\begin{equation}
 \mathscr A_p(s_c,s_b)\ge c\mathscr L_\delta^2,
 \qquad
 \mathscr A_p(S_p,s_b)\ge c\mathscr L_\delta^2.
 \label{eq:anchor-toolkit-core-action}
\end{equation}
Hence every fixed algebraic or \(e^{C S_\delta}\) loss times
\(e^{-c\mathscr A_p}\) is \(O(\delta^P)\) for every fixed \(P\).

There is a uniformly bounded \(C^2\) current frame
\begin{equation}
 T_N(r)^{-1}\mathcal J_N^{\rm fv}(r)T_N(r)
 =\begin{pmatrix}S_N^s(r)&0\\0&\lambda_{c,N}(r)\end{pmatrix},
 \qquad \|e^{S_N^s(r)t}\|\le Ce^{-D\gamma t/2}.
 \label{eq:anchor-toolkit-core-frame}
\end{equation}
In this frame the exact quotient equation is
\begin{equation}
 u'=\delta^{-1}\mathcal J_N^{\rm fv}(r_p(s))u
       +\mathcal R_p(s)u_s+f(s).
 \label{eq:anchor-toolkit-core-equation}
\end{equation}
Its coefficient and frame perturbation number satisfies
\begin{equation}
 \eta_{\rm core}(\delta)
 \le C\left\{
  \frac{\delta\mathfrak b_{\delta,\varkappa}}{S_\delta}
 +\frac{\delta^{1-2\vartheta_*}}{S_\delta}
 +\frac{\delta^{1-2\vartheta_*}}{S_\delta^2}
 +\frac1{S_\delta^2}+\frac\delta{S_\delta}\right\}=o(1).
 \label{eq:anchor-toolkit-core-perturbation}
\end{equation}

At \(s_b\), let
\begin{equation}
 \mathscr K_b=\{\phi:\ell_b^u\phi(0)=0\},\qquad
 \Gamma_b\phi=\phi-\jmath_b^u\ell_b^u\phi(0),
 \qquad (\jmath_b^ua)(\theta)=e_b^ua.
 \label{eq:anchor-toolkit-entry-coordinate}
\end{equation}
Let \(\ell_T\) be a bounded terminal row whose pairing with the transported
unstable solution \(h^u\), normalized at entry, satisfies
\begin{equation}
 |\ell_T h_T^u|\ge\mu_*>0.
 \label{eq:anchor-toolkit-terminal-margin}
\end{equation}
With the source norm
\begin{equation}
 \|f\|_{\mathcal F_p}
 =\operatorname*{ess\,sup}_{s\in I_p}
       \frac\delta{a_p(s)}\|f(s)\|_N,
 \label{eq:anchor-toolkit-source-norm}
\end{equation}
the mixed operator
\begin{equation}
 \mathbb L_pu=
 \left(u'-\delta^{-1}\mathcal J_N^{\rm fv}(r_p)u
             -\mathcal R_pu_\bullet,
       \Gamma_bu_{s_b},\ell_Tu_{S_p}\right)
 \label{eq:anchor-toolkit-mixed-operator}
\end{equation}
is an isomorphism, uniformly in \(N,\delta,p\), and
\begin{equation}
 \sup_{s\in I_p}\|u_s\|_{p,s,\sharp}
 \le C\left(\|f\|_{\mathcal F_p}
       +\|\Gamma_bu_{s_b}\|_{p,s_b,\sharp}
       +|\ell_Tu_{S_p}|\right).
 \label{eq:anchor-toolkit-mixed-inverse}
\end{equation}
The homogeneous process has a codimension-one stable bundle and a
one-dimensional unstable bundle, with
\begin{align}
 \|\mathcal V_p(s,t)\Pi_p^s(t)\|
 &\le Ce^{-c\mathscr A_p(s,t)},\notag\\
 \|\{\mathcal V_p(s,t)|_{E_p^u(t)}\}^{-1}\|
 &\le Ce^{-c\mathscr A_p(s,t)},\qquad s\ge t.
 \label{eq:anchor-toolkit-core-dichotomy}
\end{align}

The inverse is stable under the nonlinear fixed-clock problems used to
construct the intrinsic sheet.  More precisely, use the action-balanced
domain and range norms obtained from
\cref{eq:anchor-toolkit-core-action-norm,eq:anchor-toolkit-source-norm} and
the natural entry and terminal trace scales.  If \(\mathscr R_p(X)\) is the
exact nonlinear residual, \(\bar X_p\) its reference path, and
\(\mathbb B_p=D\mathscr R_p(\bar X_p)\), then
\begin{align}
 \|\mathbb B_p^{-1}\|_{Y_p^A\to X_p^A}&\le C,\notag\\
 \|\mathbb B_p^{-1}\{D\mathscr R_p(X)
                  -D\mathscr R_p(Y)\}\|_{X_p^A\to X_p^A}
 &\le C\|X-Y\|_{X_p^A},
 \label{eq:anchor-toolkit-balanced-two-point}
\end{align}
as long as \(X,Y\) remain in the declared compatible-history tube.
Consequently, on a ball
\(\|X-\bar X_p\|_{X_p^A}\le R\mathfrak b_{\delta,\varkappa}\),
\begin{equation}
 \sup_X\|\mathbb B_p^{-1}
       \{D\mathscr R_p(X)-\mathbb B_p\}\|
 \le CR\mathfrak b_{\delta,\varkappa}<\frac13.
 \label{eq:anchor-toolkit-balanced-newton}
\end{equation}
No flat bound on the second derivative of the stable graph is asserted.
For fixed (N), the selected mixed solution, unstable generator, and entry
conormal are \(C^1\) in \(p\) in fixed strong-to-weak charts, with uniform
first-jet bounds after the displayed action scaling.
\end{proposition}

\begin{proof}
The speed and graph estimates in
\cref{prop:anchor-toolkit-reference-hits} give
\(a_p(s)\asymp|r_p(s)|\).  On
\([\rho_{b,\delta},\rho_{c,\delta}]\), integration of the capped rate gives
\cref{eq:anchor-toolkit-core-action}; the logarithmic-square exponent
dominates every fixed power of \(\delta^{-1}\) and every
\(e^{CS_\delta}\) factor.

The frame in \cref{eq:anchor-toolkit-core-frame} is the invariant graph
frame from \cref{prop:anchor-toolkit-critical-splitting}.  The exact
resource quotient differs from its frozen diagonal by the generated graph
correction, the bounded delay layers, the Volterra term, and the derivative
of \(T_N(r_p(s))\).  In their respective current and history norms these
four contributions are bounded by the five terms on the right of
\cref{eq:anchor-toolkit-core-perturbation}.

For zero perturbation, solve the stable current equation forward and the
one-dimensional unstable equation from its terminal datum.  The history on
the first delay interval is prescribed by \(\Gamma_bu_{s_b}\), and later
histories are generated by the method of steps.  This gives a mixed Green
operator with the bounds in \cref{eq:anchor-toolkit-core-dichotomy}.  The
small number in \cref{eq:anchor-toolkit-core-perturbation} closes a Neumann
series.  Replacing the terminal current row by \(\ell_T\) is a scalar Schur
operation: if \(u_0\) has zero unstable terminal datum, then
\[
 u=u_0+h^u\frac{b-\ell_T(u_0)_T}{\ell_Th_T^u}.
\]
The margin \cref{eq:anchor-toolkit-terminal-margin} proves
\cref{eq:anchor-toolkit-mixed-inverse}.

Finally, the RFDE field is polynomial in current values and bounded linear
in each delayed value.  Compatible-history reconstruction and event
evaluation are bounded on the declared strong windows.  Taylor's formula
therefore gives \cref{eq:anchor-toolkit-balanced-two-point}; the action
weights are precisely those already used by the Green inverse, so no
additional unstable factor appears.  The Newton estimate follows on the
stated ball.  Differentiating the selected equation after solving it gives
the \(C^1\) assertion; this again avoids an operator-topology derivative of
the rough-history process.
\end{proof}
